% Springer Nature LaTeX template version 3.1 (December 2024).
% Prepared for Network Modeling Analysis in Health Informatics and Bioinformatics.
\pdfminorversion=7
\pdfsuppresswarningpagegroup=1
\documentclass[pdflatex,sn-mathphys-ay]{sn-jnl}

\usepackage{graphicx}
\usepackage{amsmath,amssymb,amsfonts}
\usepackage{amsthm}
\usepackage{booktabs}
\usepackage{array}
\usepackage{tabularx}
\usepackage{microtype}
\usepackage{xurl}

\graphicspath{{./}{../../figs/}}

\theoremstyle{definition}
\newtheorem{definition}{Definition}

\newcommand{\tas}{\textit{Arabidopsis thaliana}}
\newcommand{\ath}{\textit{A.\ thaliana}}
\newcommand{\hs}{\textit{Homo sapiens}}
\newcommand{\TAU}{\tau}
\newcommand{\xRightarrow}[1]{\stackrel{#1}{\Longrightarrow}}
\newcolumntype{Y}{>{\raggedright\arraybackslash}X}

\hypersetup{
  pdftitle={Observational comparison of qualitative biological network models: a reproducible Petri-net framework},
  pdfauthor={Carlos Ramirez Ovalle},
  pdfkeywords={biological network modelling, Petri nets, weak bisimulation, model comparison, systems bioinformatics, reproducible research}
}

\begin{document}

\title[Observational comparison of biological network models]
{Observational comparison of qualitative biological network models: a
reproducible Petri-net framework}

\author*[1]{\fnm{Carlos} \sur{Ramirez Ovalle}}\email{carlosovalle@javerianacali.edu.co}

\affil*[1]{\orgdiv{Department of Natural Sciences and Mathematics},
\orgname{Pontificia Universidad Javeriana Cali},
\orgaddress{\street{Calle 18 No. 118-250, Via a Pance}, \city{Cali},
\state{Valle del Cauca}, \country{Colombia}}}

\abstract{Models with similar biological-network structure can differ in event
order, branching and hidden steps. We present a reproducible framework that
converts labelled qualitative models to reachable transition systems and
evaluates strong and weak bisimulation, one-way simulation and trace distance
under a declared observational interface. Validation uses six synthetic pairs
with construction-level ground truth, six controls from two externally authored
GINsim models, and 284 public Boolean networks inferred for three HPN-DREAM
breast-cancer cell lines. One structure-only medoid per HPN-DREAM family was
fixed without reading held-out mTOR-inhibitor responses. Across the 12 controlled
comparisons, Python and mCRL2 202607.0 agreed on all 24 strong/weak decisions;
mCRL2 also matched all predeclared weak-trace controls. Weak bisimulation was
correct in 6/6 synthetic cases, while trace equality had one false positive and
a structural profile made three errors. The implementations also agreed on all
18 strong/weak decisions for nine data-conditioned cross-cell comparisons: six
had one-way simulation and three were not comparable. CASPOTS held-out
compatibility was reproduced, with excess RMSE from 0 to 0.0126, but native
medoids ranked first or tied in only two of three cell lines and their advantage
over cross-cell controls was inconclusive (exact $p=0.25$). Thus test-set
compatibility, not cell-specific prognostic value, is supported. An illustrative
plant--animal application remains a conditional comparison of encoded models,
not evidence of organism-level equivalence. Source files, hashes, tests,
notebook and generated results are public.}

\keywords{biological network modelling, Petri nets, weak bisimulation,
model comparison, systems bioinformatics, reproducible research}

\maketitle

\section{Introduction}\label{sec:introduction}

Cross-system studies often begin by aligning genes, proteins or interactions.
Orthology, phenologs and network alignment can reveal conserved components and
topology \citep{McGary2010,Singh2008,Quimbaya2012}, but structural correspondence
does not determine whether two executable models offer the same observable
choices in the same order. Conversely, models can differ in internal detail while
retaining the same externally visible behaviour. This distinction matters when a
network is used to support a cross-species or model-organism argument.

Petri nets provide an established qualitative representation of biological
regulation, causality, conflict and concurrency \citep{Murata1989,Koch2005,
ChaouiyaPetriBio2007,Heiner2008}. Executable-cell-biology approaches similarly
emphasise that a mechanistic claim should have an explicit operational semantics
\citep{Regev2002,Fisher2007}. Once a Petri net is expanded into its reachable
labelled transition system, behavioural relations from model checking can ask
whether each move of one model can be matched by the other \citep{Milner1989,
BaierKatoen2008,Sangiorgi2011}. These relations examine event order and branching
rather than only graph summaries.

This article contributes a reproducible framework that combines these ideas into
an auditable model-comparison workflow. Its primary contribution is computational:
given two explicit qualitative models and an observational interface, it returns
strong or weak equivalence, a simulation relation, or non-comparability, together
with a trace distance and diagnostic outputs. The framework is evaluated first
on synthetic pairs whose expected relations are fixed by construction. Every
transition system is independently checked with the mCRL2 toolset, reducing
dependence on a single implementation. The complete import-to-comparison path is
then exercised on two externally authored GINsim models. A structural-profile
baseline and a trace-only baseline expose what information is lost when order or
branching is ignored. Runtime is measured as model size and the candidate
relation grow. Finally, a blinded HPN-DREAM analysis connects the formal classes
to unseen perturbation-response data and tests whether cell-line-specific model
families retain specificity outside their learning data \citep{Hill2016,
Razzaq2018}.

An Arabidopsis--animal comparison of cancer-relevant modules is retained as an
illustrative application. It was motivated by comparative-genomics studies
\citep{Quimbaya2012,ClavijoBuritica2023} and the cancer-hallmark literature
\citep{Hanahan2011,Hanahan2022}. Importantly, the application does not validate
\tas{} as a model of human cancer. It reports relations between hand-curated
qualitative models under declared interfaces and identifies model-dependent
hypotheses that would require independent biological testing.

The six contributions are:
\begin{enumerate}
\item an explicit input--output procedure for behavioural comparison of
qualitative biological network models;
\item a synthetic benchmark with strong, weak, partial and negative ground-truth
cases, cross-checked against the independent mCRL2 implementation;
\item external-model controls generated from public mammalian cell-cycle and
p53--Mdm2 logical networks, with source hashes and initial conditions recorded;
\item blinded comparison of three independently inferred HPN-DREAM Boolean-model
families against held-out mTOR-inhibitor phosphoproteomic responses;
\item transparent structure-only and trace-only baselines plus runtime scaling;
and
\item an end-to-end reproducibility package containing code, tests, public and
curated models, an executed notebook and generated outputs.
\end{enumerate}

The operational sequence is summarised in Fig.~\ref{fig:pipeline}. Biological
delimitation and interface declaration precede formal encoding, reachability
construction and the graded decision rule. This separation makes the modelling
assumptions visible before any equivalence calculation is performed.

\begin{figure*}[t]
\centering
\includegraphics[width=0.96\textwidth]{Fig1.pdf}
\caption{Six-stage workflow. Stages 1--3 delimit the model-comparison question
and observational interface; stages 4--6 construct executable Petri nets,
compare their reachable behaviour and assign a model-level relation}
\label{fig:pipeline}
\end{figure*}

\section{Methods}\label{sec:methods}

\subsection{Problem statement and scope}\label{sec:problem}

The input is a pair $(N_1,N_2)$ of labelled Petri-net models and an explicitly
declared set $\Sigma$ of observable labels. The output is a relation between the
\emph{models}: strong equivalence, weak observational equivalence, one-way or
mutual simulation, or non-comparability. No step of the procedure infers an
organism-level equivalence. Biological interpretation is conditional on the
model construction, initial conditions and selected interface.

\subsection{Labelled Petri nets and reachable behaviour}\label{sec:petri}

A labelled place/transition net is
\[
N=(P,T,\mathrm{pre},\mathrm{post},\ell,m_0),
\]
where $P$ is the set of places, $T$ the set of transitions,
$\mathrm{pre}$ and $\mathrm{post}$ are input and output multisets,
$\ell:T\rightarrow\Sigma\cup\{\TAU\}$ assigns an observable label or the
internal label $\TAU$, and $m_0$ is the initial marking. Standard interleaving
semantics are used: an enabled transition fires by consuming its input tokens and
producing its output tokens. Reachability enumeration yields a labelled
transition system $\mathcal{L}(N)=(S,\Sigma\cup\{\TAU\},\rightarrow,s_0)$.

The implementation enumerates reachable markings explicitly. All curated case-
study nets are safe, so their tokens act as Boolean activation flags. A configurable
state cap stops enumeration if the chosen abstraction causes state explosion.

\subsection{Behavioural relations}\label{sec:relations}

For an observable $a\in\Sigma$, a weak move
$s\xRightarrow{a}s'$ abbreviates zero or more internal steps, one $a$-labelled
step, and zero or more internal steps. Weak bisimulation requires every observable
or internal move on either side to be matchable by the other side while preserving
the relation between successor states \citep{Milner1989,Sangiorgi2011}. Strong
bisimulation treats $\TAU$ as an ordinary label. A one-way weak simulation requires
the matching condition in one direction only and therefore detects model
containment without asserting equivalence.

The implementation computes greatest fixed points by starting with all
$|S_1||S_2|$ state pairs and repeatedly removing pairs that violate the transfer
condition. This simple implementation prioritises auditability over asymptotic
optimisation. Its candidate relation requires $O(|S_1||S_2|)$ storage; runtime
also depends on outgoing transitions and the number of refinement passes.

As a quantitative linear-time complement, the prefix-closed weak observable
language $L_k$ contains traces through depth $k$. The reported distance is
\begin{equation}
d_k(\mathcal{L}_1,\mathcal{L}_2)
=1-\frac{|L_k(\mathcal{L}_1)\cap L_k(\mathcal{L}_2)|}
{|L_k(\mathcal{L}_1)\cup L_k(\mathcal{L}_2)|}.
\label{eq:distance}
\end{equation}
Weak bisimilarity implies weak trace equality, but the converse is false.
Consequently, $d_k$ is used for ranking and diagnostics, never as a replacement
for the branching-time relation.

\subsection{Graded decision rule}\label{sec:decision}

The decision rule is fixed before a pair is evaluated:
\begin{enumerate}
\item report \emph{strong equivalence} when strong bisimulation holds;
\item otherwise report \emph{weak equivalence} when weak bisimulation holds;
\item otherwise report \emph{mutual simulation} or the direction of one-way
simulation; and
\item report \emph{not comparable} when neither simulation direction holds.
\end{enumerate}
The trace distance is attached to, rather than substituted for, this categorical
result.

\subsection{Synthetic benchmark with construction-level ground truth}
\label{sec:benchmark-method}

Six deterministic cases were specified independently of the biological models.
A 12-step labelled workflow is used for the first five cases. The identity case
copies it exactly. Silent refinement subdivides selected observable edges with
internal $\TAU$ steps. Label mismatch replaces a recurring observable. Label
order swap preserves topology and label counts while changing two consecutive
events. Added branch retains the reference workflow and adds one observable
outcome. The sixth case is the standard branching-time distinction
$a.(b+c)$ versus $a.b+a.c$: the trace sets agree, but one model resolves its
choice before the other. Table~\ref{tab:synthetic} records the relation expected
from each construction.

\begin{table*}[t]
\caption{Synthetic benchmark and predeclared model-level relation. The observed
class was computed without using the expected label}\label{tab:synthetic}
\centering
\small
\begin{tabularx}{\textwidth}{@{}>{\raggedright\arraybackslash}p{0.18\textwidth}>{\raggedright\arraybackslash}p{0.18\textwidth}YY@{}}
\toprule
Case & Controlled change & Expected class & Observed class \\
\midrule
Identity & exact copy & strong equivalence & strong equivalence \\
Silent refinement & insert internal steps & weak equivalence & weak equivalence \\
Label mismatch & replace one observable family & not comparable & not comparable \\
Label order swap & preserve counts, change order & not comparable & not comparable \\
Added branch & preserve workflow, add outcome & reference simulated by candidate & reference simulated by candidate \\
Trace-equivalent branching & move point of choice & candidate simulated by reference & candidate simulated by reference \\
\botrule
\end{tabularx}
\end{table*}

\subsection{Independent mCRL2 oracle}\label{sec:mcrl2-method}

To test implementation independence, each benchmark LTS was exported to the
Aldebaran AUT interchange format and evaluated with \texttt{ltscompare} from
mCRL2 202607.0 \citep{Bunte2019}. The external oracle computed strong
bisimilarity, weak bisimilarity and exact weak-trace equivalence, treating
\texttt{tau} as internal. These are separate executions of mCRL2 algorithms on
the exported systems; the Python verdict is not supplied as input. Agreement is
reported for strong and weak bisimulation, while exact weak-trace results are
compared with the trace relation fixed by each construction. The current
\texttt{ltscompare} interface does not provide a weak simulation preorder, so
one-way weak-simulation classifications are not claimed to have an mCRL2 oracle.

\subsection{Controls from public logical models}\label{sec:public-method}

Two externally authored models were obtained from the public GINsim repository
\citep{Naldi2018}: the Boolean mammalian cell-cycle model of Traynard et al.
\citep{Traynard2016}, distributed as SBML-qual \citep{Chaouiya2013SBML}, and the
multilevel p53--Mdm2 DNA-damage model associated with Abou-Jaoud\'e et al.
\citep{AbouJaoude2009}, distributed as GINML. The exact downloaded files, source
URLs and SHA-256 digests are retained in the repository. For the cell-cycle
model, \texttt{CycD=1} and all other components were initially zero. For the
p53--Mdm2 model, the stored GINsim state \texttt{damage} was used.

Both rule sets were interpreted with unitary asynchronous semantics: in one
transition, one enabled non-constant component moves one level toward its
logical target. The observable label records the component and direction
(\texttt{component\_up} or \texttt{component\_down}). Breadth-first enumeration
then constructs the reachable LTS. For each public base model, three controls
were predeclared: an exact copy must be strongly and weakly equivalent; inserting
five internal steps must break strong but preserve weak and weak-trace
equivalence; and renaming one reachable update family must break all three
relations. These controls test parsing, state generation, AUT export and formal
comparison on external dynamics. They do not establish that the published
logical rules are complete biological representations.

\subsection{Blinded HPN-DREAM validation with held-out data}
\label{sec:hpn-method}

The data-backed validation uses the public HPN-DREAM phosphoproteomic time series
for BT20, BT549 and MCF7 breast-cancer cell lines \citep{Hill2016} and the
corresponding CASPOTS Boolean-network families \citep{Razzaq2018}. The archived
families contain 72, 191 and 21 verified
networks, respectively. Files were recovered from public commit
\texttt{95e3c74} of the CASPOTS repository; the prior-knowledge network, learning
sets, model families and mTOR-inhibitor test sets are retained with SHA-256
digests. Learning and test files remain separate throughout the workflow.

One representative per cell line was fixed before reading any test response.
For each family, pairwise Jaccard similarity was computed between active Boolean
clauses. The model with maximum mean within-family similarity was selected; exact
ties were resolved lexicographically on the active-clause tuple using rational
arithmetic. This structure-only rule selected zero-based rows 24, 137 and 9 for
BT20, BT549 and MCF7. An anti-leakage unit test raises an exception if medoid
selection attempts to open a test file.

The observational interface was the intersection of seven measured
PI3K/MAPK/mTOR readouts: 4EBP1-pS65, AKT-pT308, BAD-pS112, MAPK-pT202/Y204,
MEK1-pS217/S221, mTOR-pS2448 and p70S6K-pT389. The perturbation set was fixed as
the exact intersection of the three test panels: mTOR inhibition alone,
IGF1 plus mTOR inhibition and 4EBP1 plus mTOR inhibition. Test values determined
neither model selection nor interface membership. Each selected Boolean model
was interpreted with unitary asynchronous semantics. Measured time-zero values
initialized relevant readouts, other relevant nodes were zero, stimuli were
clamped active and inhibitors inactive. Interface updates were labelled by
component and direction; all other updates were $\TAU$. The resulting nine
cross-cell pair--condition comparisons were evaluated by both implementations.

Experimental distance was the RMSE between the two cell-line profiles over
common post-zero times and interface readouts. Time-zero measurements initialized
the LTS but were excluded from this outcome to prevent data reuse. Formal classes were ordered from weak
bisimulation through mutual and one-way simulation to non-comparability. Their
association with experimental distance was evaluated with Spearman's $\rho$ and
an exact one-sided test that permuted the three pair labels within each condition
($3!^3=216$ permutations). Truncated trace distance was a secondary metric.

Held-out compatibility was also recomputed with CASPOTS/Clingo. The network was
fixed while the solver minimized the published weighted threshold-repair
objective. Because streamed optimization models need not be final, the runner
consumed the solve handle to completion and added deterministic secondary
objectives for continuous squared error and source-row index. Clingo certified
the final optimum, and each of 12 scores was reproduced twice identically. Each
blind medoid was scored on all three test sets; selecting the best network from a
whole family on its own test set was reported separately as a retrospective
oracle, not a blind estimate. Native specificity was tested by exact sign
permutation of the three native-versus-mean-cross differences. Setting every
unmeasured relevant initial node to one was retained as a state-space stress test;
its largest LTS had 27,648 states and was not entered into a quadratic exact
cross-model relation.

\subsection{Non-formal baselines}\label{sec:baselines}

Two transparent baselines isolate the information contributed by the formal
relation. The structural-profile score averages four similarities: state count,
edge count, observable-label multiset and out-degree multiset. Each count
similarity is the smaller-to-larger ratio; multiset similarity is weighted
Jaccard. A score of at least 0.9 is classified as equivalent. This feature-only
baseline deliberately has no event-order semantics and is not presented as a
state-of-the-art network aligner. The trace baseline declares equivalence when
$d_8=0$ and therefore has order but no branching-time discrimination.

Binary equivalence ground truth is positive only for the identity and silent-
refinement cases. Strong equivalence is included within weak equivalence for this
binary evaluation; simulations are not counted as equivalence.

\subsection{Scalability experiment}\label{sec:scalability-method}

Linear workflows of 8, 16, 32, 64, 96 and 128 transitions were compared with an
identical copy and a silently refined version. For each pair, strong bisimulation,
weak bisimulation, both simulation directions and $d_8$ were measured separately.
Wall-clock times are medians of three repetitions. The reference run used Python
3.11.8 on an Apple M3 Pro with 18~GiB memory. Runtime values are machine-dependent;
state counts, edge counts, candidate relation sizes and classifications are
deterministic.

\subsection{Illustrative biological model set}\label{sec:case-method}

The application contains five pairs of compact, safe Petri nets: DNA repair and
genome instability (GIM), energy metabolism (DCE), cell-cycle control (SPS), cell
death/autophagy (RCD), and an immune module (AID). Components were selected from
published plant--human comparative studies \citep{Quimbaya2012,
ClavijoBuritica2023}; dedicated reviews informed the qualitative representation
of plant DNA-damage response, cell-cycle control and autophagy
\citep{Yoshiyama2013,DeVeylder2007,Marshall2019}.

Places represent conditions or molecular states, transitions represent curated
events, and the initial marking represents the declared starting context. A
transition is observable only when its label belongs to the cross-system
interface; species-specific intermediate steps are represented as $\TAU$.
Transition-level provenance is stored in the repository and checked for complete
name coverage. These nets are manually curated qualitative abstractions. They do
not encode kinetics, concentrations, stoichiometric uncertainty or tissue
context, and they were not inferred from independent experimental data.

Figure~\ref{fig:gim-petri} exposes the complete GIM encodings rather than only
their computed verdict. Both nets implement damage sensing, checkpoint control,
repair and restoration. Their terminal mechanisms differ: the animal/human net
contains apoptosis, whereas the plant net contains an internal SMR step followed
by differentiation/endoreduplication. The common label \texttt{exit\_cycle} is a
declared coarse observation, not evidence that those terminal mechanisms are
biologically identical.

\begin{figure*}[t]
\centering
\begin{minipage}[t]{\textwidth}
  \centering
  \includegraphics[width=0.98\textwidth]{Fig2a.pdf}
  \par\smallskip\textbf{(a) Animal/human GIM model}
\end{minipage}
\par\medskip
\begin{minipage}[t]{\textwidth}
  \centering
  \includegraphics[width=0.98\textwidth]{Fig2b.pdf}
  \par\smallskip\textbf{(b) Arabidopsis GIM model}
\end{minipage}
\caption{Curated DNA-damage Petri nets used in the illustrative comparison:
(a) animal/human and (b) Arabidopsis. Circles denote places and the initial
token; rounded boxes denote transitions.
Observable labels appear in brackets, while grey transitions labelled $\tau$
are internal. The diagrams are generated from the same executable definitions
used by the analysis}
\label{fig:gim-petri}
\end{figure*}

The interface is treated as a modelling hypothesis. Diagnostics include targeted
single-label changes, random label permutations, seed sensitivity, a trace-depth
sweep and silent-refinement invariance. The permutation values quantify
sensitivity to scrambled interfaces inside these models; they are not biological
significance tests and do not provide external validation.

\subsection{Software and reproducibility}\label{sec:repro-method}

The formal engine, public-model parsers and benchmark use the Python standard
library. Plotting and tabulation use NumPy, pandas, Matplotlib and NetworkX;
Petri-net artwork uses pydot and Graphviz. Independent checking uses the
hash-pinned mCRL2 202607.0 release. Submission figures are exported as vector PDF with embedded
fonts, and colour is supplemented by labels, symbols or line styles. Fixed seeds
control all synthetic modifications and interface permutations. The repository contains
the exact model registry, benchmark generator, unit tests, executed Jupyter
notebook, environment specifications, continuous-integration workflow and script
that regenerates every reported table and figure. Runtime CSV files are excluded
from bit-for-bit checks because wall-clock measurements are hardware-dependent;
all other generated results are deterministic.



\section{Results}\label{sec:results}

\subsection{The synthetic suite recovers every predeclared relation}
\label{sec:benchmark-results}

The formal classification matched the construction-level class in all six cases
(Table~\ref{tab:synthetic}). Identity was strongly bisimilar. Inserting internal
steps broke strong equivalence but preserved weak equivalence. Label mismatch and
label-order change were both rejected. The added branch produced the intended
one-way simulation, while the trace-equivalent branching pair was also separated
by simulation direction rather than incorrectly promoted to equivalence.

For the binary weak-equivalence task, weak bisimulation achieved 6/6 correct
decisions. Trace equality achieved 5/6 because it produced a false positive on
the branching-time case. The structural profile achieved 3/6: it produced false
positives for label mismatch and label-order swap and a false negative for silent
refinement (Table~\ref{tab:baseline}; Fig.~\ref{fig:benchmark}).

\begin{table}[t]
\caption{Binary equivalence accuracy on six construction-level cases}
\label{tab:baseline}
\centering
\begin{tabular}{lccc}
\toprule
Method & Accuracy & False positive & False negative \\
\midrule
Weak bisimulation & 1.000 & 0 & 0 \\
Trace equality ($k=8$) & 0.833 & 1 & 0 \\
Structural profile ($\geq0.9$) & 0.500 & 2 & 1 \\
\botrule
\end{tabular}
\end{table}

\begin{figure*}[t]
\centering
\includegraphics[width=\textwidth]{Fig3.pdf}
\caption{Independent benchmark. Left: correctness and predicted binary class for
each case and method. Right: accuracy against the equivalence labels fixed by
construction. The trace baseline misses the branching-time distinction; the
structural profile additionally misses event-order and silent-refinement cases}
\label{fig:benchmark}
\end{figure*}

\subsection{Independent software and public-model controls agree}
\label{sec:external-results}

The public cell-cycle model contained 13 Boolean components and generated 826 of
8,192 possible states with 3,413 reachable transitions. The four-component
multilevel p53--Mdm2 model generated 17 of 36 possible states and 26 transitions
(Table~\ref{tab:public-models}). Thus, the external controls exercise the method
on substantially larger and differently encoded state spaces than the compact
curated application.

\begin{table*}[t]
\caption{Public logical models used for external workflow validation. Full space
is the product of component level counts; only reachable states enter the formal
comparison}\label{tab:public-models}
\centering
\scriptsize
\setlength{\tabcolsep}{2pt}
\begin{tabularx}{\textwidth}{@{}>{\raggedright\arraybackslash}p{2.25cm}lrrrr>{\raggedright\arraybackslash}X@{}}
\toprule
Model & Format & Comp. & Full space & Reachable & Edges & Initial context \\
\midrule
Mammalian cell cycle & SBML-qual & 13 & 8,192 & 826 & 3,413 & \texttt{CycD=1}; others zero \\
p53--Mdm2 DNA damage & GINML & 4 & 36 & 17 & 26 & stored state \texttt{damage} \\
\botrule
\end{tabularx}
\end{table*}

All six public-model controls recovered their predeclared relations in both
implementations (Table~\ref{tab:public-controls}). Silent refinement added five
states to each candidate without changing weak behaviour. The observable
perturbations were rejected by strong bisimulation, weak bisimulation and exact
weak-trace equivalence. Across these six controls and the six synthetic pairs,
Python and mCRL2 agreed on 12/12 strong and 12/12 weak decisions. The mCRL2 exact
weak-trace result also matched all 12 predeclared trace controls, including
acceptance of the trace-equivalent branching pair that both bisimulation
implementations reject.

\begin{table*}[t]
\caption{Public-model controls. Entries list strong/weak/weak-trace outcomes;
Python reports the first two relations and mCRL2 reports all three. Y=yes and
N=no}\label{tab:public-controls}
\centering
\scriptsize
\setlength{\tabcolsep}{3pt}
\begin{tabular*}{\textwidth}{@{\extracolsep{\fill}}llccccc}
\toprule
Model & Control & \shortstack{States\\ref./cand.} & \shortstack{Edges\\ref./cand.} & \shortstack{Expected\\S/W/T} & \shortstack{Python\\S/W} & \shortstack{mCRL2\\S/W/T} \\
\midrule
Cell cycle & exact copy & 826/826 & 3,413/3,413 & Y/Y/Y & Y/Y & Y/Y/Y \\
Cell cycle & silent refinement & 826/831 & 3,413/3,418 & N/Y/Y & N/Y & N/Y/Y \\
Cell cycle & observable perturbation & 826/826 & 3,413/3,413 & N/N/N & N/N & N/N/N \\
p53--Mdm2 & exact copy & 17/17 & 26/26 & Y/Y/Y & Y/Y & Y/Y/Y \\
p53--Mdm2 & silent refinement & 17/22 & 26/31 & N/Y/Y & N/Y & N/Y/Y \\
p53--Mdm2 & observable perturbation & 17/17 & 26/26 & N/N/N & N/N & N/N/N \\
\botrule
\end{tabular*}
\end{table*}

\subsection{Held-out data support compatibility but not prognostic specificity}
\label{sec:hpn-results}

The structure-only representatives contained 23--29 active clauses
(Table~\ref{tab:hpn-medoids}). Across the three shared perturbations, their
reachable LTSs contained 2--44 states. None of the nine cross-cell comparisons
was weakly bisimilar: six showed one-way simulation and three were not comparable.
Python and mCRL2 agreed on every strong and weak decision (18/18), and mCRL2
rejected weak-trace equivalence in all nine pairs.

\begin{table}[t]
\caption{HPN-DREAM model families and blinded structure-only representatives}
\label{tab:hpn-medoids}
\centering
\small
\begin{tabular}{lrrrr}
\toprule
Cell line & Family & Medoid row$^a$ & Clauses & Mean Jaccard \\
\midrule
BT20  & 72  & 24  & 29 & 0.815 \\
BT549 & 191 & 137 & 25 & 0.854 \\
MCF7  & 21  & 9   & 23 & 0.851 \\
\bottomrule
\end{tabular}
\par\smallskip
\raggedright\footnotesize $^a$ Zero-based row; selected from model structure only.
\end{table}

The primary ordinal class had a small positive association with between-cell
experimental distance ($\rho=0.091$), but the exact condition-stratified test was
inconclusive ($p=0.111$). Trace distance was negatively associated with the data
distance ($\rho=-0.596$, $p=0.556$), showing that the truncated linear-time index
is not an empirical similarity surrogate in this small panel.

CASPOTS reproduced compatibility with the held-out test data. Native medoid RMSE
was 0.33024 for BT20, 0.31335 for BT549 and 0.27726 for MCF7; corresponding
discretization RMSE was 0.32944, 0.30079 and 0.27726. The native medoid was best
or tied for BT20 and MCF7, but not for BT549 (Fig.~\ref{fig:hpn-validation}a).
Across the three targets, the native excess-RMSE advantage over the mean of both
cross-cell medoids was 0.00212 and was not significant under the exact paired
sign test ($p=0.25$). The independently selected formal class likewise did not
significantly order experimental distances (Fig.~\ref{fig:hpn-validation}b).
These results validate held-out compatibility and deterministic reproduction of
the published scoring logic, while directly rejecting a claim of demonstrated
cell-specific predictive discrimination.

\subsection{Runtime grows with the candidate relation}\label{sec:scaling-results}

Across the two scaling series, complete runtime rose from approximately
0.4~ms for 9-state reference models to approximately 0.54~s for the largest
pairs. The 128-step identity pair contained $129\times129=16{,}641$ candidate
state pairs; silent refinement increased the candidate to 145 states and the
relation to 18,705 pairs. Both series retained their expected classifications at
every size (Fig.~\ref{fig:scaling}). Two simulation computations accounted for
the largest runtime share. These measurements establish feasibility for compact
and medium qualitative modules, not for genome-scale state spaces.

\begin{figure*}[t]
\centering
\includegraphics[width=0.96\textwidth]{Fig4.pdf}
\caption{Scalability on deterministic workflows. Left: median total runtime
against candidate state-pair count on logarithmic axes. Right: decomposition of
the largest run into strong bisimulation, weak bisimulation, two simulations and
trace distance. Timings are machine-dependent; model sizes and outcomes are
deterministic}
\label{fig:scaling}
\end{figure*}

\subsection{Case-study outputs are conditional relations between curated models}
\label{sec:case-results}

The ten curated Petri nets generated reachable systems of 5--14 states. Under the
declared interfaces, the GIM, DCE and SPS model pairs were weakly, but not
strongly, bisimilar and had $d_6=0$. The RCD plant model was simulated by the
animal model but not conversely ($d_6=0.143$). The AID pair satisfied neither
simulation direction ($d_6=0.429$; Table~\ref{tab:case}).

\begin{table*}[t]
\caption{Relations between the curated case-study model pairs. These are
model-level, interface-conditional outputs; $P$ denotes the plant model and
$A$ the animal model}\label{tab:case}
\centering
\small
\begin{tabular*}{\textwidth}{@{\extracolsep\fill}llccccc}
\toprule
Module & Encoded process & Strong & Weak & $P\preceq A$ & $A\preceq P$ & $d_6$ \\
\midrule
GIM & DNA repair & no & yes & yes & yes & 0.000 \\
DCE & energy metabolism & no & yes & yes & yes & 0.000 \\
SPS & cell-cycle control & no & yes & yes & yes & 0.000 \\
RCD & cell death/autophagy & no & no & yes & no & 0.143 \\
AID & immune control & no & no & no & no & 0.429 \\
\botrule
\end{tabular*}
\end{table*}

In GIM, the animal terminal fate of apoptosis and the plant terminal fate of
differentiation/endoreduplication are represented through the common coarse
observable \texttt{exit\_cycle}. The weak-equivalence result therefore says that
the two encoded models agree at this selected resolution. A finer interface that
distinguishes those outcomes would change the classification. This example makes
the conditional nature of the method explicit rather than treating the interface
as a neutral measurement device.

The corresponding reachability systems in Fig.~\ref{fig:gim-lts} make the weak
relation inspectable. Observable paths for damage, checkpoint, repair,
restoration and cycle exit can be matched, while additional internal paths are
absorbed by $\TAU$ closure. The figure therefore connects the Petri-net
construction to the state-pair relation used by the algorithm.

\begin{figure*}[t]
\centering
\begin{minipage}[t]{0.49\textwidth}
  \centering
  \includegraphics[width=\textwidth]{Fig5a.pdf}
  \par\smallskip\textbf{(a) Animal/human reachability system}
\end{minipage}\hfill
\begin{minipage}[t]{0.49\textwidth}
  \centering
  \includegraphics[width=\textwidth]{Fig5b.pdf}
  \par\smallskip\textbf{(b) Arabidopsis reachability system}
\end{minipage}
\caption{Reachability systems generated from the GIM Petri nets: (a)
animal/human and (b) Arabidopsis. Numbered nodes are reachable markings, the
initial node is distinguished by its fill, solid
labelled arcs are observable moves and dashed arcs are internal $\tau$ moves}
\label{fig:gim-lts}
\end{figure*}

The RCD pair supplies a contrasting one-way result (Fig.~\ref{fig:rcd-petri}).
Both encodings include stress, mTOR inhibition, autophagy and survival. The
animal model additionally exposes canonical caspase apoptosis, whereas the plant
model proceeds through an internal metacaspase event. Consequently, the plant
behaviour is simulated by the animal model under the declared interface, but the
reverse direction fails. This is a containment statement about the two encodings,
not a claim that either net exhausts cell-death biology.

\begin{figure*}[t]
\centering
\begin{minipage}[t]{\textwidth}
  \centering
  \includegraphics[width=0.94\textwidth]{Fig6a.pdf}
  \par\smallskip\textbf{(a) Animal RCD model with an observable caspase-apoptosis branch}
\end{minipage}
\par\medskip
\begin{minipage}[t]{\textwidth}
  \centering
  \includegraphics[width=0.94\textwidth]{Fig6b.pdf}
  \par\smallskip\textbf{(b) Arabidopsis RCD model with an internal metacaspase branch}
\end{minipage}
\caption{Curated cell-death/autophagy Petri nets illustrating the source of the
one-way simulation result: (a) animal and (b) Arabidopsis. Symbols and transition styling follow
Fig.~\ref{fig:gim-petri}}
\label{fig:rcd-petri}
\end{figure*}

Randomly subdividing internal structure preserved every categorical result, as
required by weak semantics. Renaming individual observables changed the expected
model relations, confirming that interface labels affect the computed result.
Label-permutation diagnostics separated the three weakly equivalent curated
pairs from scrambled versions, but these tests only assess the internal
specificity of the selected interface. They cannot establish biological validity
because the original labels and networks were manually curated
(Fig.~\ref{fig:diagnostics}).

\begin{figure*}[t]
\centering
\begin{minipage}[t]{0.47\textwidth}
  \centering
  \includegraphics[width=\textwidth]{Fig7a.pdf}
  \par\smallskip\textbf{(a) Silent-refinement invariance}
\end{minipage}\hfill
\begin{minipage}[t]{0.51\textwidth}
  \centering
  \includegraphics[width=\textwidth]{Fig7b.pdf}
  \par\smallskip\textbf{(b) Scrambled-interface reference}
\end{minipage}
\caption{Internal robustness diagnostics for the curated case study: (a)
verdict preservation under 300 silent refinements per module and (b) observed
distance against 2,000 scrambled-interface distances. Bars and numeric labels
report refinement invariance; violin distributions show the label-permutation
reference and points show observed distances. Reported
$p$- and Benjamini--Hochberg $q$-values quantify model-internal interface
specificity and are not experimental biological significance tests}
\label{fig:diagnostics}
\end{figure*}

\begin{figure*}[t]
\centering
\includegraphics[width=\textwidth]{Fig8.pdf}
\caption{Blinded HPN-DREAM validation. (a) CASPOTS excess RMSE above the
discretization reference for every structure-only source medoid and held-out
target cell line; blue boxes mark native source--target combinations. (b)
Between-cell profile RMSE under the primary formal class for three shared
perturbations. Horizontal lines are class medians. The exact test permutes cell
pairs within perturbation; neither panel supports robust native-cell
specificity}
\label{fig:hpn-validation}
\end{figure*}

\section{Discussion}\label{sec:discussion}

\subsection{What the framework establishes}\label{sec:discussion-adds}

The framework turns a qualitative network comparison into a reproducible
question with explicit inputs and falsifiable model-level outputs. The synthetic
benchmark shows why graph summaries and trace sets are not interchangeable with
branching-time relations. A topology- and label-count-preserving order change can
remain structurally similar while losing simulation; conversely, two systems can
share all finite traces in the benchmark yet resolve a choice differently. Weak
bisimulation detects both distinctions while tolerating silent implementation
detail.

The synthetic benchmark is small by design but independent of the biological
application. mCRL2 removes dependence on the in-house fixed-point implementation,
while the two repository models show that the parser and comparison pipeline
operate on externally authored logical dynamics. Together these layers remove
the circular claim that success on hand-built case-study models alone validates
the software. The HPN-DREAM analysis adds a distinct evidential layer: model
selection is blind to held-out values, the formal outcomes are independently
checked and the same model representatives are confronted with experimental
profiles. The case study instead demonstrates provenance and conditional
interpretation on a domain-motivated model registry.

\subsection{What the framework does not establish}\label{sec:discussion-limits}

The case-study outputs do not show that plant and human molecular activity is
experimentally equivalent. The models omit kinetics, quantities, noise,
stoichiometry, cell type and environmental context. Their small state spaces make
exact checking and visual audit possible but restrict biological resolution. The
observational interface can also merge mechanistically different outcomes, as
the GIM \texttt{exit\_cycle} example demonstrates.

The biological registry is grounded primarily in two comparative studies and
selected reviews. Different curators could choose different arcs, initial
markings or observables. Transition-name coverage proves documentation
completeness, not evidential independence. Likewise, label-permutation values
measure sensitivity within the constructed models and must not be interpreted as
experimental $p$-values.

The implementation uses explicit state enumeration and a transparent fixed-point
algorithm. State explosion and the quadratic candidate relation bound limit its
use on large composed networks. The scaling experiment covers linear systems up
to 145 states and does not establish genome-scale performance. Partition-
refinement algorithms, symbolic state representations and compositional checking
are natural extensions.

The transformed GINsim controls still have a precise limit: their expected
relations arise from exact copy, silent refinement and label perturbation. The
HPN-DREAM experiment addresses that weakness with distinct public model families
and held-out measurements, but it does not yield a positive validation of
prognostic specificity. Only three cell lines had both verified families and
usable test sets, the primary association had nine pair--condition observations,
and the native advantage was not significant. CASPOTS RMSE is a minimum-distance
compatibility score over admissible trajectories, not a unique time-series
prediction \citep{Razzaq2018}. The negative specificity result therefore prevents
using this analysis to claim that behavioural class predicts response distance.

Finally, the structural baseline is intentionally simple. It tests whether basic
graph features suffice on controlled counterexamples; it is not a comprehensive
comparison with every network-alignment method. Future benchmarks should include
more repository models, prospectively expert-labelled positive and negative
cross-model relations, biological replicates, alternative alignment algorithms
and blinded interface specifications.

\subsection{Practical use}\label{sec:discussion-use}

The method is most appropriate when researchers already have two explicit
qualitative models and need to document the exact resolution at which they agree
or diverge. A defensible workflow should predeclare the interface, retain fine
outcome distinctions when they matter to the research question, report plausible
alternative interfaces and validate any biological interpretation against data
not used to construct the models. In that role, the framework is a model-auditing
and hypothesis-prioritisation tool rather than a substitute for biological
experimentation.

\section{Conclusion}\label{sec:conclusion}

We presented a reproducible framework for comparing qualitative biological
network models through strong and weak bisimulation, simulation and trace
distance. Synthetic constructions, an independent mCRL2 oracle and two public
GINsim models test complementary parts of the software path. The blinded
HPN-DREAM analysis further shows that independently inferred models can be
checked against unseen intervention data: formal decisions were reproducible and
the models retained low held-out compatibility error, but neither class--data
concordance nor native-cell specificity was established. Complete agreement on
controlled relations supports implementation correctness; the negative data
result defines the present inferential boundary. The Arabidopsis--animal
application demonstrates the workflow but is deliberately interpreted at model
level: its results depend on curated networks and declared interfaces and do not
prove organism-level functional equivalence.

\backmatter

\section*{Statements and Declarations}

\noindent\textbf{Funding} The author received no specific funding for this work.

\noindent\textbf{Competing interests} The author declares no financial or non-financial
interests directly or indirectly related to this work.

\noindent\textbf{Ethics approval and consent to participate} Not applicable. The study
uses computational models and previously published aggregate information; no
human participants, human samples or live animals were involved.

\noindent\textbf{Consent for publication} Not applicable.

\noindent\textbf{Data availability} Curated definitions, hash-verified public GINsim
models, HPN-DREAM/CASPOTS model families and normalized perturbation-response
files, deterministic benchmark metadata and all generated result tables are available at
\url{https://github.com/cero1979/BisimulationMol}. No confidential or participant-
level data were used.

\noindent\textbf{Materials availability} The complete reproducibility package, including
environment specifications and the executed notebook, is available in the same
public repository.

\noindent\textbf{Code availability} Source code is released under the MIT License at
\url{https://github.com/cero1979/BisimulationMol}. The commands
\texttt{make external-validation}, \texttt{make hpn-validation},
\texttt{make figures}, \texttt{make test}
and \texttt{make notebook} regenerate the analysis, validate expected relations
with mCRL2 and execute the notebook.

\noindent\textbf{Generative AI statement} OpenAI Codex was used during revision
for programming assistance, reproducibility checks and language editing. The
author inspected the source artifacts, executed the analyses, verified all
reported outputs and assumes full responsibility for the manuscript.

\noindent\textbf{Author contributions} Carlos Ramirez Ovalle: conceptualization,
methodology, software, formal analysis, validation and visualization; writing---
original draft; writing---review and editing.



\bibliography{references}

\end{document}
